% Working draft, 18 November 2026. Prepared for discussion; not a submission.
% Sources: handwritten outline (2 Nov 2026), docs/theory_reference/single_scenario_support_pedagogical.tex,
% SUPPORT_LOSS_EXPERIMENTS_V1.md, experiments/support_scaling_20261026/reports/RESULTS_20261026_1434_LIVE.json.
\documentclass[11pt,a4paper]{article}
\usepackage[margin=24mm]{geometry}
\usepackage[T1]{fontenc}
\usepackage{lmodern,microtype}
\usepackage{amsmath,amssymb,amsthm,mathtools}
\usepackage{booktabs,enumitem,xcolor}
\usepackage[hidelinks]{hyperref}
\setlist{nosep,leftmargin=*,topsep=3pt}
\newtheorem{theorem}{Theorem}
\newtheorem{proposition}[theorem]{Proposition}
\theoremstyle{remark}
\newtheorem{remark}[theorem]{Remark}
\newcommand{\R}{\mathbb R}
\newcommand{\E}{\mathbb E}
\newcommand{\Z}{\mathcal Z}
\newcommand{\hh}{\hat h}
\newcommand{\Ls}{L_\sigma}
\DeclareMathOperator*{\argmin}{arg\,min}
\newcommand{\todo}[1]{\textcolor{red!70!black}{[#1]}}

\title{Decision-focused learning for linear recourse problems with random right-hand side:\\
a convex support surrogate}
\author{Jonathan Hornewall\\[2pt]{\normalsize Working draft for discussion, 18 November 2026}}
\date{}

\begin{document}
\maketitle

\begin{abstract}
\noindent
Decision-focused learning (DFL) is well developed for contextual linear
programs, where the unknown parameter is the cost vector. Contextual
\emph{linear recourse} problems, where a second-stage decision is taken after
the uncertainty is revealed, are a richer and practically more common family,
but they lack a counterpart to the SPO+ loss. For two-stage problems with a
random right-hand side (demand), we propose the \emph{support surrogate}: a loss
on a predicted demand vector that (i) upper-bounds the true decision cost,
(ii) is convex in the prediction when the first stage is continuous, (iii) is
the optimal value of a single LP whose dual solution is a subgradient, with no
smoothing and no hyperparameter, and (iv) is consistent in a strong sense: \emph{every}
minimiser of the expected surrogate induces a decision that is optimal for the
true conditional stochastic program. On four small benchmark families with 100
training pairs, the method has the lowest mean cost gap among eleven methods in
all four; on a partial larger-scale campaign it is lowest in four of seven
benchmarks and level with the best in a fifth. With an integer first stage, on the
benchmark of the SFGE paper, it reduces test cost by 11--35\% relative to SFGE. This note states the idea, the main
guarantee with its short proof, and the current empirical evidence. It is a
draft meant to convey the idea, not a finished manuscript.
\end{abstract}

\section{Motivation}

Optimization problems often have parameters that are unknown at the time the
decision has to be made. Frequently, machine learning is used to produce
statistically accurate predictions of these parameters. Such predictions carry
no general guarantee of producing high-quality decisions, which may require
taking the full distribution of possible outcomes into account. This is the
principal motivation for decision-focused learning, which trains prediction
models to minimise decision regret rather than statistical deviation.

Most work has focused on contextual linear programs: linear programs with an
unknown cost vector. Contextual linear recourse problems add a second-stage
decision taken after the unknown parameter is revealed, and allow uncertainty in
the right-hand side and the constraint matrix as well. Attempts to extend DFL
to this family date from early in the field's development, but two obstacles
remain. A single predicted scenario cannot in general be the conditional mean
and still give the right decision, so it is unclear what should be predicted;
and the map from prediction to decision cost is piecewise constant or
non-convex, so it is unclear how to train.

Demand (right-hand-side) uncertainty is the most important special case:
it is the most common form in applications, there is a lot of activity on it in
applied fields, and it is where the foundations are least developed. This note
treats that case.

\paragraph{What is proposed.}
A single convex surrogate for the decision cost of a demand prediction,
obtained from the dual of the second stage, together with
\begin{itemize}
\item an upper bound on the true decision cost and a consistency theorem
      stronger than Fisher consistency (Section~\ref{sec:surrogate});
\item a single-LP form of the loss, with exact subgradients from its dual and no hyperparameter (Section~\ref{sec:training});
\item experiments against ten baselines (Section~\ref{sec:empirics});
\item the extension to integer first-stage decisions, with experiments on the SFGE paper's benchmark (Section~\ref{sec:integer}).
\end{itemize}

\paragraph{Closest prior art (to be written up properly).}
Decision rules and predict-then-optimise approaches for contextual stochastic
programming; SPO+ for cost vectors; Lex-SPO, SFGE and decision-focused scenario
learning for recourse problems; the CO-augmented ML line and the survey of
Bengio et al. The methods we compare against empirically are listed in
Section~\ref{sec:empirics}. \todo{references and a careful positioning to add}

\section{Setting}

Given a context $x$, we choose a first-stage decision $z\in\Z$, then observe a
random demand vector $H$ and pay a linear recourse cost:
\[
 \min_{z\in\Z}\ c^\top z+\E\bigl[Q(z,H)\mid x\bigr],
 \qquad
 Q(z,h)=\min_{y\ge0}\{q^\top y:\ Wy=h-Tz\}.
\]
Here $c,T,W,q$ are known and fixed; only the demand $H$ is random, and $h$ denotes a realisation of it. By LP duality the
recourse cost is the support function of the second-stage dual feasible set,
\[
 Q(z,h)=\sigma(h-Tz),\qquad
 \sigma(r)=\sup_{\lambda\in\Delta}\lambda^\top r,\qquad
 \Delta=\{\lambda:\ W^\top\lambda\le q\}.
\]
Write $G(z,h)=c^\top z+\sigma(h-Tz)$ for the total cost of decision $z$ under
demand $h$, and $F_x(z)=\E[G(z,H)\mid x]$ for the conditional expected cost,
with minimum $F_x^\star$ attained at some $z^\star$.

\paragraph{Predict, then decide.}
A model produces a single predicted demand $\hh=\hh(x)$. The decision is
obtained by solving the \emph{single-scenario} problem
\[
 \hat z(\hh)\in\argmin_{z\in\Z}G(z,\hh),\qquad V(\hh)=\min_{z\in\Z}G(z,\hh).
\]
The quantity we care about is the realised cost of that decision,
$G(\hat z(\hh),h)$; its regret is the difference to the best achievable cost.
As a function of $\hh$ this is non-convex and its gradient is uninformative
wherever the decision is locally constant.

\section{The support surrogate}\label{sec:surrogate}

A support function is sublinear: $\sigma(a+b)\le\sigma(a)+\sigma(b)$. Splitting
the residual $h-T\hat z(\hh)=(\hh-T\hat z(\hh))+(h-\hh)$ gives
\begin{align*}
 G(\hat z(\hh),h)
 &=c^\top \hat z(\hh)+\sigma\bigl(\hh-T\hat z(\hh)+h-\hh\bigr)\\
 &\le c^\top \hat z(\hh)+\sigma\bigl(\hh-T\hat z(\hh)\bigr)+\sigma(h-\hh)
 \;=\;V(\hh)+\sigma(h-\hh).
\end{align*}
We define the \textbf{support surrogate}
\begin{equation}\label{eq:loss}
 \boxed{\ \Ls(\hh,h)=V(\hh)+\sigma(h-\hh)
        =\min_{z\in\Z}\bigl\{c^\top z+\sigma(\hh-Tz)\bigr\}+\sigma(h-\hh).\ }
\end{equation}
The first term is the cost of planning for the prediction; the second is the
recourse cost of correcting the prediction to the realisation, priced by the
problem's own second stage.

\begin{proposition}[Upper bound]\label{prop:bound}
For every prediction $\hh$, every realisation $h$ and every optimal
single-scenario decision $\hat z(\hh)$,
$\ G(\hat z(\hh),h)\le\Ls(\hh,h)$, with equality at $\hh=h$, where both equal
the perfect-information cost $V(h)$.
\end{proposition}

So $\Ls(\hh,h)-V(h)\ge0$ upper-bounds the regret against perfect information and
vanishes for a perfect prediction. In expectation, however, the surrogate is
\emph{not} minimised by the conditional mean, and that is the point.

\begin{theorem}[Minimisers of the surrogate are decision-optimal]\label{thm:main}
Fix a context $x$ and let $z^\star\in\argmin_{z}F_x(z)$. Then
\[
 \min_{\hh}\ \E\bigl[\Ls(\hh,H)\mid x\bigr]=F_x^\star,
 \quad\text{attained at }\hh^\star=Tz^\star,
\]
and for every $\hh$ and every optimal single-scenario decision $\hat z(\hh)$,
\[
 0\ \le\ F_x\bigl(\hat z(\hh)\bigr)-F_x^\star\ \le\
 \E\bigl[\Ls(\hh,H)\mid x\bigr]-F_x^\star .
\]
In particular, any prediction that minimises the expected surrogate has zero
expected decision regret.
\end{theorem}

\begin{proof}
Taking expectations in Proposition~\ref{prop:bound},
$\E[\Ls(\hh,H)\mid x]\ge F_x(\hat z(\hh))\ge F_x^\star$ for all $\hh$. For the
reverse inequality set $\hh^\star=Tz^\star$. Since $z^\star$ is feasible in the
minimisation defining $V$ and $\sigma(0)=0$,
\[
 \Ls(\hh^\star,h)=V(Tz^\star)+\sigma(h-Tz^\star)
 \le c^\top z^\star+\sigma(Tz^\star-Tz^\star)+\sigma(h-Tz^\star)=G(z^\star,h).
\]
Taking expectations gives $\E[\Ls(\hh^\star,H)\mid x]\le F_x^\star$.
\end{proof}

\begin{remark}[Stronger than Fisher consistency]
The theorem does not only say that the surrogate minimiser minimises regret.
It gives a calibration inequality with constant one: surrogate excess risk
bounds decision regret for \emph{every} prediction, so approximate minimisation
of the surrogate yields approximately optimal decisions. A predictor that
minimises the surrogate therefore yields an optimal decision rule, which
suggests DFL as an attractive alternative to learning decision rules directly.
\end{remark}

\begin{remark}[No tie-breaking assumption]
Nothing is assumed about how the solver selects among multiple optima. Any
decision that is optimal for a surrogate-minimising prediction is optimal for
the true conditional problem. For LPs, where predictions routinely induce
several optima, this is a useful guarantee.
\end{remark}

\begin{remark}[What is being predicted]
The optimal report $\hh^\star=Tz^\star$ is the demand that the optimal decision
serves exactly. It is generally neither the conditional mean nor a possible
realisation of $H$. For the newsvendor it is the critical quantile. The
surrogate is what tells a regression model to aim for this point: for the
newsvendor $\Ls$ reduces to the pinball loss, and in general it plays the role
that SPO+ plays for cost vectors.
\end{remark}

\begin{remark}[Scope and caveats]
The statements require the losses to be finite on the set of admissible
predictions: $\sigma(h-\hh)<\infty$ means that the second stage can absorb the
prediction error in either direction (e.g.\ shortage and surplus variables).
Without complete recourse the report must be restricted to a domain where
this holds. The same argument extends to a random recourse cost $q$ (decode
with an upper bound or with the conditional mean of $q$) and, under a
containment condition on the dual sets, to a random $W$; these are not used in
the experiments below.
\end{remark}

\section{Training}\label{sec:training}

\begin{proposition}[Convexity and subgradient]
If $\Z$ is a polyhedron, $\Ls(\cdot,h)$ is convex and piecewise linear. Let
$\lambda_0$ be an optimal dual multiplier of the constraint $Wy=\hh-Tz$ in the
single-scenario LP defining $V(\hh)$, and $\lambda_h$ an optimal dual
solution of the residual LP $\min_{y\ge0}\{q^\top y:Wy=h-\hh\}$. Then
\[
 \lambda_0-\lambda_h\in\partial_{\hh}\Ls(\hh,h).
\]
\end{proposition}

\paragraph{The loss as a single LP.}
Both terms of the loss are minimisations, so they combine into one problem:
\begin{equation}\label{eq:singlelp}
 \Ls(\hh,h)=\min_{z\in\Z,\ y_0\ge0,\ y_h\ge0}
 \bigl\{\,c^\top z+q^\top y_0+q^\top y_h\ :\ Tz+Wy_0=\hh,\ \ Wy_h=h-\hh\,\bigr\}.
\end{equation}
This is the two-stage problem with two recourse blocks: $y_0$ serves the
predicted demand and $y_h$ pays for the residual $h-\hh$. The prediction enters
only through the right-hand side, so $\Ls(\cdot,h)$ is the value function of an
LP in its right-hand side, which is convex and piecewise linear, and the
proposition is LP sensitivity: if $(\lambda_0,\lambda_h)$ are optimal
multipliers of the two constraint blocks, the derivative of the right-hand side
with respect to $\hh$ is $(I,-I)$, which gives $\lambda_0-\lambda_h$. One solve
of~\eqref{eq:singlelp} per sample therefore returns the loss, its subgradient
and the decision. Because the two blocks share no variables, the solve also
splits into the single-scenario problem and a small recourse LP, which is how
we implement it.

Training a predictor $\hh_w$ on pairs $(x_i,h_i)$ is then ordinary stochastic
subgradient descent on $\frac1N\sum_i\Ls(\hh_w(x_i),h_i)$:
\begin{enumerate}
\item predict $\hh_i=\hh_w(x_i)$;
\item solve~\eqref{eq:singlelp} for $(\hh_i,h_i)$; keep its value and the
      multipliers $\lambda_{0,i}$, $\lambda_{h,i}$;
\item backpropagate $\lambda_{0,i}-\lambda_{h,i}$ through $\hh_w$.
\end{enumerate}
At deployment only the single-scenario problem is solved. Compared with the
alternatives we tested, the loss needs one LP per sample, of the size of a two-scenario problem, no differentiation through a solver, no smoothing or
regularisation of the optimisation problem, no sampling, and no loss-specific
hyperparameter. Convexity is in the prediction, not in neural-network weights.

Two further facts, stated without proof here:
\begin{itemize}
\item \emph{Exact fitting.} For an affine predictor on fixed features, the
      empirical surrogate minimisation is itself a single LP (a MILP if the
      first stage is integer), so a global empirical minimiser is available.
\item \emph{Error decomposition.} Expected regret of the deployed policy is
      bounded by the predictor-class approximation error of the surrogate, plus
      twice the uniform deviation between empirical and population surrogate
      risk, plus the training and solver gaps. The first term vanishes when
      the class can represent $x\mapsto Tz^\star(x)$.
\end{itemize}

\section{Empirical evidence}\label{sec:empirics}

\subsection{Protocol}

All benchmarks are two-stage LPs with continuous decisions and random demand
only; $W$ and $q$ are the native recourse data of each model. Each fit sees
\textbf{100 training pairs} $(x_i,h_i)$, one observed demand per context, and 32
validation pairs. Predictors are affine or a two-hidden-layer MLP (width 64),
trained with Adam for 300 epochs. Hyperparameters and architecture are chosen on
one development seed using validation cost only, then frozen for ten
independent confirmation seeds. Learners use only the supplied training data
and the known optimisation structure.

The metric is the cost gap of the deployed decision relative to a reference
obtained by solving a 100-scenario conditional SAA at each of 100 test
contexts, as a percentage: per seed,
$100\times\mathrm{mean}_i(C_i-V_i)/\mathrm{mean}_i|V_i|$.

\paragraph{Baselines.}
Non-contextual SAA; $k$NN SAA; ER-SAA (empirical-residual SAA); MSE regression
followed by the single-scenario problem; a differentiable quadratic-regularised
solver layer trained on decision cost; SFGE (score-function gradient
estimation, authors' code); Lex-SPO (authors' formulation); a static
contextual scenario-generation method of Islip et al.\ (authors' code); and two
direct policies that output a first-stage decision, one projected at inference
and one trained through a differentiable projection.
\todo{exact citations}

\subsection{Four small benchmark families (complete)}

The families are resource procurement, a shipment network, a two-layer network
inventory problem and capacity assignment: 12--37 first-stage variables, 12--26
recourse variables, six demand and six context coordinates, with a nonlinear
context-dependent demand law. Table~\ref{tab:v1} reports all eleven methods.

\begin{table}[h]
\centering\small
\caption{Cost gap (\%), mean $\pm$ standard error over ten confirmation seeds,
$N=100$ training pairs. Architecture selected per method on the development
seed (A affine, M MLP, S not applicable). Bold: lowest mean.}\label{tab:v1}
\begin{tabular}{@{}lrrrr@{}}
\toprule
Method & Resource proc. & Shipment netw. & Netw. inventory & Capacity assign.\\
\midrule
Non-contextual SAA & $61.30\pm1.31$ & $48.93\pm2.34$ & $80.07\pm1.69$ & $79.15\pm0.79$\\
$k$NN SAA & $22.84\pm0.82$ & $27.34\pm1.76$ & $36.70\pm0.94$ & $46.43\pm2.20$\\
ER-SAA & $5.29\pm1.00$ (M) & $13.40\pm1.96$ (M) & $13.69\pm0.76$ (M) & $11.66\pm0.60$ (M)\\
MSE + single scenario & $5.31\pm0.45$ (M) & $25.62\pm1.45$ (M) & $19.50\pm1.21$ (M) & $16.33\pm0.82$ (M)\\
Quadratic solver layer & $5.42\pm0.30$ (A) & $10.58\pm1.08$ (M) & $17.91\pm0.88$ (M) & $11.81\pm0.35$ (M)\\
SFGE & $5.67\pm0.35$ (A) & $11.71\pm0.97$ (M) & $16.48\pm0.78$ (M) & $15.27\pm0.59$ (M)\\
Lex-SPO & $7.22\pm0.27$ (A) & $11.05\pm1.23$ (A) & $18.73\pm0.38$ (A) & $17.06\pm0.51$ (A)\\
Islip et al.\ (static) & $4.95\pm0.29$ (M) & $24.59\pm4.73$ (M) & $17.53\pm0.42$ (A) & $11.50\pm0.37$ (M)\\
Affine policy, projected & $8.98\pm0.39$ (A) & $12.08\pm1.17$ (A) & $19.17\pm0.38$ (A) & $17.40\pm0.46$ (A)\\
Direct policy, diff.\ proj. & $5.15\pm0.29$ (A) & $11.03\pm1.08$ (M) & $13.74\pm0.65$ (M) & $13.12\pm0.52$ (M)\\
\midrule
\textbf{Support surrogate} & $\mathbf{2.39\pm0.09}$ (A) & $\mathbf{10.00\pm1.01}$ (M) & $\mathbf{11.69\pm0.77}$ (M) & $\mathbf{10.73\pm0.36}$ (M)\\
\bottomrule
\end{tabular}
\end{table}

\paragraph{Reading the table.}
The support surrogate has the lowest mean gap in all four families. On
resource procurement the gap is roughly halved relative to the best competitor
($2.39$ vs.\ $4.95$). On the other three it leads the strongest competitor,
which differs by family (the quadratic layer, ER-SAA and Islip et al.\
respectively), by between half a point and two points.

The improvement over MSE regression with the same architecture and the same
single-scenario decoder ($2.4$ vs.\ $5.3$, $10.0$ vs.\ $25.6$, $11.7$ vs.\
$19.5$, $10.7$ vs.\ $16.3$) isolates the effect of the loss: the model class
and the deployed optimisation problem are identical, only the training target
differs.

\subsection{Larger instances (partial campaign)}

A scaling campaign over problem size and training-set size was launched on
26 October and is \textbf{incomplete}: the compute window closed before most
cells finished, and continuation on the CERMICS cluster is planned. What exists
is the first cell: the ``large'' instances (e.g.\ 64 resources, 96 demand
classes, 32 suppliers; 50-product newsvendor) and the two original benchmarks
of our scenario-generation paper, all at $N=100$. Lex-SPO, the quadratic layer
and the direct policies are missing for most benchmarks in this snapshot, so
Table~\ref{tab:scale} is a comparison against a reduced roster.

\begin{table}[h]
\centering\small
\caption{Cost gap (\%), mean $\pm$ standard error over ten seeds, large
instances and original benchmarks, $N=100$. Snapshot of 26 October; partial
campaign. Architecture selected on the development seed. Bold: lowest mean
among the methods shown.}\label{tab:scale}
\setlength{\tabcolsep}{2.5pt}\footnotesize
\begin{tabular}{@{}lrrrrrrr@{}}
\toprule
 & Res. & Shipm. & Netw. & Capac. & News- & Res.\ alloc. & Shipment\\
Method & proc. & netw. & invent. & assign. & vendor & (orig.) & (orig.)\\
\midrule
Non-contextual SAA & $68.7$ & $54.0$ & $78.8$ & $82.4$ & $348.7$ & $42.0$ & $61.5$\\
$k$NN SAA & $34.8$ & $33.0$ & $40.6$ & $38.7$ & $154.7$ & $7.8$ & $19.1$\\
MSE + single scenario & $7.4$ & $27.8$ & $15.8$ & $15.5$ & $65.2$ & $2.1$ & $12.7$\\
SFGE & $21.3$ & $58.9$ & $85.1$ & $91.7$ & $276.0$ & $7.0$ & $11.8$\\
ER-SAA & $\mathbf{4.00\pm0.32}$ & $7.04\pm0.49$ & $11.38\pm0.47$ & $10.50\pm0.29$ & $49.7\pm2.1$ & $\mathbf{0.75\pm0.12}$ & $3.12\pm0.94$\\
Islip et al.\ (static) & $4.96\pm0.19$ & $6.74\pm0.58$ & $10.36\pm0.38$ & $9.56\pm0.22$ & $41.5\pm1.6$ & $1.03\pm0.14$ & $2.72\pm0.40$\\
\midrule
\textbf{Support surrogate} & $5.31\pm0.21$ & $\mathbf{5.55\pm0.65}$ & $\mathbf{9.96\pm0.41}$ & $\mathbf{9.23\pm0.16}$ & $41.9\pm2.4$ & $1.47\pm0.23$ & $\mathbf{1.94\pm0.40}$\\
\bottomrule
\end{tabular}
\end{table}

\paragraph{Reading the table.}
\begin{itemize}
\item The surrogate has the lowest mean in four of seven benchmarks (both
      shipment problems, network inventory, capacity assignment).
\item On the newsvendor it is level with Islip et al.\ and with the quadratic
      layer ($41.3$, one of the two benchmarks where that method has
      results); the percentages are large because the reference cost of this
      benchmark is small.
\item ER-SAA is ahead on resource procurement and on the original resource
      allocation problem, by $1.3$ and $0.7$ points.
\item SFGE degrades sharply at this size, and MSE regression remains far
      behind in the shipment problems.
\end{itemize}
Lex-SPO has results only on the two original benchmarks ($8.6$, on eight
seeds, and $16.4$).

\section{Integer first-stage decisions}\label{sec:integer}

Proposition~\ref{prop:bound} and Theorem~\ref{thm:main} use only the
sublinearity of the continuous second stage. They make no use of convexity of
$\Z$, so they hold verbatim when the first stage is integer or mixed-integer:
the surrogate still upper-bounds the decision cost, $\hh^\star=Tz^\star$ still
minimises its expectation, and every decoder optimum at a minimiser is
optimal. The single-problem form~\eqref{eq:singlelp} also carries over, as a MILP. What
is lost is convexity of $V$ in $\hh$: $V$ is now a minimum over integer
decisions of convex functions, and an LP multiplier is no longer a global
subgradient.

\paragraph{Differentiating the loss.}
In the integer case we train a Gaussian predictor, as SFGE does:
$\hh\sim\pi_w(\cdot\mid x)=\mathcal N(\mu_w(x),\mathrm{diag}(s^2))$ with a
learned variance $s^2$, and we minimise the expected loss
$\E_{\hh\sim\pi_w}[\Ls(\hh,h)]$. The decision deployed at test time is
$\hat z(\mu_w(x))$. The two terms of the loss are differentiated separately.
\begin{itemize}
\item \emph{Residual term, by duality.} $\sigma(h-\hh)$ is a convex function of
      $\hh$ with subgradient $-\lambda_h$, exactly as in the continuous case.
      Writing $\hh=\mu_w(x)+s\,\varepsilon$ with $\varepsilon\sim\mathcal N(0,I)$,
      this subgradient is backpropagated through the sample.
\item \emph{Decoder term, by the score-function identity.} $V(\hh)$ is the
      value of a MILP. Its expectation is differentiated without
      differentiating the MILP:
      \[
       \nabla_w\,\E_{\hh\sim\pi_w}\bigl[V(\hh)\bigr]
       =\E_{\hh\sim\pi_w}\Bigl[\bigl(V(\hh)-V(\mu_w(x))\bigr)\,
         \nabla_w\log\pi_w(\hh\mid x)\Bigr],
      \]
      where the value at the mean, $V(\mu_w(x))$, is a baseline that reduces
      variance. This needs only optimal values of the single-scenario MILP.
\end{itemize}
We call this estimator \emph{score}. In a second estimator, \emph{dual}, the
decoder term is also differentiated by duality: with the integer decision fixed
at $\hat z(\hh)$, the second stage is an LP in $\hh$, and its multiplier
$\lambda_0$ is the gradient of $V$ on the piece selected by $\hat z(\hh)$; it
is backpropagated through the same sample. SFGE applies the score-function
identity to the whole decision regret. Here the identity is applied only to
$V$, and the term that contains the observed demand has an exact gradient.

\paragraph{Benchmark.}
We use the weighted set multi-cover problem of the SFGE paper, with its
generator and its 800/100/100 train/validation/test split. The first stage
chooses integer cover counts $z$; unmet demand is penalised:
\[
 \min_{z\in\mathbb Z_+^{n}}\ c^\top z+\E\bigl[p^\top(H-Az)_+\mid x\bigr],
\]
so $\sigma(r)=p^\top r_+$ and the surrogate is
$\Ls(\hh,h)=V(\hh)+p^\top(h-\hh)_+$. There are two sizes ($5$ demand
coordinates and $25$ integer variables; $10$ and $50$) and three shortage-penalty
multipliers, each with five generated instances and three splits: 90
configurations in total. Predictors are linear or MLPs with one to three hidden
layers of width 64, selected per method on validation cost.

\begin{table}[h]
\centering\small
\caption{Weighted set multi-cover, integer first stage. Mean test cost
(thousands) and change in test cost of the support surrogate relative to SFGE,
averaged over paired configurations. Architecture selected on validation for
each method. Bold: lowest mean cost.}\label{tab:wsmc}
\begin{tabular}{@{}lrrrrrr@{}}
\toprule
 & & & \multicolumn{2}{c}{Support, cost} & \multicolumn{2}{c}{Support vs.\ SFGE}\\
\cmidrule(lr){4-5}\cmidrule(l){6-7}
Size / penalty & SFGE & $k$NN SAA & score & dual & score & dual\\
\midrule
$5\times25$ / 1  & 28.2 & 29.2 & 25.6 & \textbf{25.4} & $-11.1\%$ & $-11.1\%$\\
$5\times25$ / 5  & 36.7 & 43.4 & \textbf{30.8} & 31.5 & $-17.3\%$ & $-15.3\%$\\
$5\times25$ / 10 & 45.6 & 50.6 & \textbf{33.9} & 35.7 & $-24.8\%$ & $-21.2\%$\\
$10\times50$ / 1  & 74.2 & 70.4 & 61.0 & \textbf{60.8} & $-18.3\%$ & $-18.5\%$\\
$10\times50$ / 5  & 107.4 & 102.5 & 75.1 & \textbf{75.0} & $-29.5\%$ & $-29.7\%$\\
$10\times50$ / 10 & 127.9 & 119.5 & \textbf{82.4} & 83.0 & $-35.5\%$ & $-34.6\%$\\
\bottomrule
\end{tabular}
\end{table}

With the score estimator the surrogate has lower test cost than SFGE in 88 of
the 90 configurations, and with the dual estimator in 86; both are lower than
$k$NN SAA in 89 of 90. The reductions relative to SFGE are 11--35\%, and
14--31\% relative to $k$NN SAA. The advantage grows with the shortage penalty
and with problem size, and the two estimators give similar costs. The result
depends on the predictor: with the linear predictor used in the SFGE paper the
surrogate is 5--19\% behind SFGE, while with one hidden layer or more both
estimators are ahead in all six settings at every depth. Validation never
selects the linear model for the surrogate.

\end{document}
